\section{Lessons learned}

% \subsection{Interface-Driven Modularity}
% interface 기반 motiondag로 노드 단위 개발/재구성이 가능해짐
% 기능 격리와 스왑 비용이 감소함. 인터페이스 과세분화와 이벤트 연쇄 업데이트 등과 함께 개발단계에서 있었던 썰 풀기

% \subsection{Standardized Pose I/O \& vendor conversion}
% 입력을 단일 월드좌표 6dof 스트림으로 정규화하고 sdk 변환 규칙을 잘 처리함.
% 멀티 벤더 통합 용이.

% \subsection{Zero-Copy kinematic buffer}
% 연속 메모리 버퍼 등
% 버퍼 생애주기, 관리 가이드

% \subsection{SDK-independent retargeting}
% 여긴... 해부학적으로 리타게팅을 sdk에서 분리함.
% 뭘써야함 도덕책

% \subsection{Evaluation/Sparse Input Pipelike design}
% 입력 조건에서의 사용성이나 부하(userstudy)지표 제시 및 개선점
% 뭐 다른 태스크를 시킨다면.. 이런거

% \subsection{limitation, threats to validity}
% 내가 나를 까는 입장에서 하드웨어 커버리지, 샘플 규모, 엔진에 대한 종속성..?
% 음.... 외적 타당성을 가지는 정도의 한계 작성. discussion가 겹치나?
